\scanpage{103}

\uline{$\overline{\mathrm{VII}}$}\quad\uline{Structures \unsure{triangulées|triangulaires} congruentielles.}

Une \uline{structure triangulaire} est pour nous un couple

\[
\begin{aligned}
 (S,F)\qquad&S\text{ un ensemble (``\uline{sommets}'')}\\
 &\underset{\text{ens. des \uline{faces}}}{F}\subset\unsure{\mathscr P|\mathcal P|\mathfrak P}_3(S)
\end{aligned}\tag*{(1)}
\]

satisfaisant l'axiome

\[
\deleted{\text{(T1)}}\quad\text{\textcircled{\scriptsize T1}}\qquad\deleted{\unsure{S}}\,S=\bigcup_{f\in F}f
\qquad\underset{\deleted{\text{\unsure{bien|lien}}}}{\text{donc}}\quad S\ne\varnothing\iff F\ne\varnothing.\tag*{(2)}
\]

Les structures qui nous intéressent satisfont à des axiomes supplémentaires, inspirés par le cas des $(S_M,F_M)$ où $M$ est un module libre de rang $2$ sur un $\mathbb Z/n\mathbb Z$, $n\geq3$. Posons déjà

\[
\deleted{\text{\unsure{(T1 bis)}}}\qquad\underset{\substack{\text{ens des}\\\text{\uline{arêtes}}}}{A}=\{a\in\mathscr P_2(S)\mid\exists f\in F,\ f\supset a\},\tag*{(3)}
\]

\deleted{\unsure{on}} demande

\[
\begin{aligned}
 (\mathrm{T2})\qquad &\forall a\in A,\ \text{il existe exactement deux}\\
 &f\in F\ \text{telles que }\ f\supset a
\end{aligned}\tag*{(4)}
\]

(toute arête incidente à exactement deux faces).

En d'autres termes, le triple

\[
\begin{aligned}
 \Gamma=\bigl(F_{\deleted{\unsure{1}}},A_{\deleted{\unsure{2}}},{}&
 R_{2}^{\deleted{\unsure{2}}\prime}\subset F\times S\bigr)\\
 &\shortparallel\\[-2pt]
 &\{(f,a)\in F\times A\mid f\supset a\}
\end{aligned}\tag*{(5)}
\]

est une structure de \uline{graphe strict}, ayant
% archive sheet number: 102
